\documentclass[../main.tex]{subfiles}

\begin{document}
\subsection{Antes de la fórmula integral de Cauchy}
\classdate{2026-04-17}

\begin{lemma}
    Sean \(\gamma \colon [a, b] \subseteq \mathbb{R} \to \mathbb{C}\) una cuerva cerrada
    suave por pedazos y \(z_{0} \in \mathbb{C}\setminus \gamma\). Entonces

    \begin{equation*}
        \int_{\gamma}^{}\dfrac{1}{z - z_{0}}\odif{z} = 2k\pi i,\quad \forall
        k\in \mathbb{Z}.
    \end{equation*}
\end{lemma}

\begin{proof}
    Sean \(F \colon [a, b] \subseteq \mathbb{R} \to \mathbb{C}\) dada por

    \begin{equation*}
        F(x) = \int_{a}^{x}\dfrac{\gamma^{\prime}(t)}{\gamma(t) - z_{0}}\odif{t}.
    \end{equation*}

    Como \(\gamma\) es suave por pedazos y \(z_{0}\notin \gamma([a, b])\), la
    función \(t
    \mapsto \tfrac{\gamma^{\prime}(t)}{\gamma(t) - z_{0}}\) es continua por
    pedazos en \([a,
    b]\), de modo que \(f\) está bien definida y es continua en \([a, b]\). Además, por
    TFCC, es de clase \(C^{1}\) en cada subintervalo donde \(\gamma\) lo es.

    \begin{equation*}
        F^{\prime}(x) = \dfrac{\gamma^{\prime}(x)}{\gamma(x) - z_{0}}.
    \end{equation*}

    Consideremos \(G \colon [a, b] \to \mathbb{C}\), con \(G(x) =
    \mathrm{e}^{-F(x)}\bigl(\gamma(x) - z_{0}\bigr)\) y

    \begin{equation*}
        G^{\prime}(x) = -F^{\prime}(x)\mathrm{e}^{-F(x)}\bigl(\gamma(x) - z_{0}\bigr) +
        \mathrm{e}^{-F(x)}\gamma^{\prime}(x).
    \end{equation*}

    Sustituyendo \(F^{\prime}(x)\),

    \begin{equation*}
        G^{\prime}(x) = -\mathrm{e}^{-F(x)}\dfrac{\gamma^{\prime}(x)}{\gamma(x) -
        z_{0}}\bigl(\gamma(x) - z_{0}\bigr) + \mathrm{e}^{-F(x)}\gamma^{\prime}(x) = 0.
    \end{equation*}

    Por lo que \(G\) es constante en cada subintervalo. Como \(G\) es continua en \([a,
    b]\), concluimos que \(G\) es constante en \([a, b]\) y \(\forall x\in [a, b],\;
        \mathrm{e}^{-F(x)}\bigl(\gamma(x) - z_{0}\bigr) =
        \mathrm{e}^{-F(a)}\bigl(\gamma(a)
    - z_{0}\bigr)\), pero \(F(a) = 0\), tal que

    \begin{equation*}
        \mathrm{e}^{-F(x)}\bigl(\gamma(x) - z_{0}\bigr) = \gamma(a) - z_{0}.
    \end{equation*}

    Considerando \(x = b\) y recordando que \(\gamma(b) = \gamma(a)\),

    \begin{align*}
        \mathrm{e}^{-F(b)}\bigl(\gamma(a) - z_{0}\bigr) &= \gamma(a) - z_{0},\\
        \mathrm{e}^{-F(b)} &= 1.
    \end{align*}

    De esta manera, \(\exists k\in \mathbb{Z}\) tal que \(F(b) = 2k\pi i\). Pero

    \begin{align*}
        F(b) &= \int_{a}^{b}\dfrac{\gamma^{\prime}(t)}{\gamma(t) - z_{0}}\odif{t} =
        \int_{\gamma}^{}\dfrac{1}{z - z_{0}}\odif{z},\\
        \therefore\; \int_{\gamma}^{}\dfrac{1}{z - z_{0}}\odif{z} &= 2k\pi
        i,\quad \text{para
        alguna}\; k\in \mathbb{Z}.
    \end{align*}
\end{proof}
\end{document}
